\chapter*{Введение} % * не проставляет номер
\addcontentsline{toc}{chapter}{Введение} % вносим в содержание

\noindent\textbf{Aктуальность исследования}

Современные телекоммуникационные системы сталкиваются с необходимостью обеспечения надежной передачи данных в условиях ограниченной пропускной способности, динамически изменяющейся топологии сети и высокой мобильности абонентов.
Особую актуальность эта задача приобретает в самоорганизующихся низкоскоростных сетях, где отсутствует централизованное управление, а связь между абонентами осуществляется через нестабильные радиоканалы.
Такие сети востребованы в ситуациях чрезвычайного характера, зонах природных катастроф, военных операциях и в удалённых регионах, где инфраструктура традиционных сетей коммуникации практически отсутствует.

\noindent\textbf{Степень научной разработанности проблемы}
Сегодня существуют различные подходы к решению задач маршрутизации и оптимизации передачи данных в мобильных сетях, однако существующие решения зачастую недостаточно эффективны в условиях ограниченной пропускной способности и высокой интенсивности движений абонентов.

Несмотря на значительный прогресс в теории построения беспроводных сетей, область оптимального проектирования децентрализованных самоорганизующихся сетей остается открытой проблемой. Большинство текущих подходов ориентированы на статичные условия или значительно упрощенные модели связи, тогда как современные технологии предъявляют повышенные требования к устойчивости и гибкости таких систем.

Исследование посвящено разработке алгоритма самоорганизации и маршрутизации, который обеспечит надежную передачу данных даже в сложных условиях нестабильных каналов связи. Результаты исследования имеют высокую прикладную ценность, так как могут применяться в системах мониторинга, экстренного реагирования, военной связи и других областях, где стабильность и доступность канала связи играют ключевую роль.

\noindent\textbf{Цель исследования}

Целью данной работы является разработка алгоритма функционирования самоорганизующейся низкоскоростной сети передачи данных, обеспечивающего:

\begin{enumerate}[1.]
	\item Децентрализованное управление – обмен информацией о топологии между абонентами без центрального сервера.

    \item Адаптивную маршрутизацию – прогнозирование перемещений абонентов и минимизацию разрывов связи.

    \item Оптимизацию энергопотребления – снижение нагрузки на передатчики при сохранении качества связи.

    \item Устойчивость к потерям данных – механизмы повторной передачи и восстановления фреймов.
\end{enumerate}

Для достижения поставленной цели в работе применяются:
\begin{itemize}
	\item Математическое моделирование – описание динамики сети и параметров радиоканала.

    \item Алгоритмическая разработка – создание механизмов самоорганизации и маршрутизации.

    \item Численное моделирование – оценка производительности алгоритма в различных сценариях.
\end{itemize}

Структура работы включает:
\begin{enumerate}[1.]
	\item Анализ особенностей современных самоорганизующихся низкоскоростных сетей и выявление основных проблем, влияющих на их работоспособность.
    \item Разработка математической модели, адекватно описывающей динамику и поведение таких сетей.
    \item Проектирование алгоритма адаптивной маршрутизации, учитывающего перемещение абонентов и уменьшающего вероятность разрыва соединений.
    \item Исследование эффективности метода путем имитационного моделирования. Экспериментальная проверка разработанного алгоритма на моделях виртуальной сети и реального оборудования.
\end{enumerate}

\noindent\textbf{Практическая значимость исследования}  

Работа позволяет решить конкретные практические задачи в области проектирования и эксплуатации самоорганизующихся низкоскоростных сетей передачи данных, улучшая их надежность и эффективность в условиях динамической топологии и нестабильного соединения.
    Результаты исследования способствуют повышению уровня готовности служб экстренного реагирования, аварийно-спасательных подразделений и военного ведомства к оперативному развертыванию сетей в труднодоступных районах или зонах чрезвычайных происшествий.
    Разработанный алгоритм адаптивной маршрутизации может использоваться непосредственно в действующих проектах по внедрению низкоскоростных сетей, предназначенных для экологического мониторинга, контроля промышленных объектов и сельского хозяйства.

Таким образом, результаты данного исследования обладают значительной практической ценностью и могут быть активно интегрированы в реальную производственную среду для решения важных инженерных задач.
Результаты работы могут быть использованы при проектировании автономных телекоммуникационных систем, работающих в условиях нестабильной связи и ограниченных ресурсов.

В результате проведённого исследования была достигнута следующая \textbf{научная новизна}:
\begin{itemize}
    \item Создан оригинальный метод адаптивной маршрутизации, позволяющий эффективно реагировать на изменения топологии сети, оперативно перестраивать пути передачи данных и сокращать задержку передачи данных.
    \item Выведен ряд формул и зависимостей, устанавливающих связь между параметрами мобильной сети и эффективностью работы алгоритмов маршрутизации, основанных на принципах самоорганизации.
    \item Осуществлён переход от качественной оценки недостатков традиционных методов передачи данных к количественному определению критериев оптимальности.
\end{itemize}
%%%%%%%%%%%%%%%%%%%%%%%%%%%%%%%%%%

% \textbf{Объект исследования} --- Самоорганизующиеся низкоскоростные сети передачи данных, характеризующиеся отсутствием центральной координации, высоким уровнем мобильности абонентов и ограниченными ресурсами энергии.

% \textbf{Предмет исследования} --- Механизмы адаптивного распределения трафика, оптимизация энергопотребления и построение надежных маршрутов передачи данных в условиях быстро изменяющейся топологии сети.

% \textbf{Цель исследования} --- Разработка эффективного алгоритма самоорганизации и адаптивной маршрутизации для низкоскоростных сетей передачи данных, обеспечивающего надежное функционирование сети в условиях нестабильного радиоокружения и высокой мобильности абонентов. 
% В соответствии с целью исследования, логически определяются следующие \textbf{задачи работы}:

% В процессе разработки и тестирования алгоритма самоорганизации и адаптивной маршрутизации выдвигаются следующие научные предположения и гипотезы, нуждающиеся в проверке и подтверждении:

%     Использование адаптивных стратегий маршрутизации существенно снижает потери данных и повышает надежность передачи в самоорганизующихся низкоскоростных сетях.
%     Предположительно, сочетание алгоритмов прогнозирования перемещений абонентов и динамического обновления маршрутов позволит сократить число разорвавшихся соединений и обеспечить своевременную доставку сообщений даже в условиях повышенной мобильности и низкой пропускной способности.
%     Оптимизированные механизмы экономии энергии позволяют снизить потребление ресурса батареи абонентов без ухудшения качества связи.
%     Предполагается, что сокращение числа избыточных отправлений и уменьшение активности передатчиков за счет грамотного планирования интервалов сканирования позволят увеличить срок службы автономного абонента, сохраняя высокий уровень доступности сети.
%     Комплексный подход к снижению потерь данных с применением методов коррекции ошибок улучшает показатели стабильности работы сети.Мы полагаем, что использование комбинированных механизмов восстановления утраченных кадров совместно с интеллектуальным перераспределением передаваемой нагрузки повысит коэффициент успешности доставки данных, снизив потребность в повторных передачах и увеличив общую эффективность сети.
%     Создание модели самоорганизации на основе теории графов и математического моделирования позволит достоверно оценивать характеристики сети и оптимизировать её архитектуру.
%     Есть основания полагать, что формализованный подход к исследованию свойств графа взаимосвязей абонентов позволит лучше понимать внутренние зависимости и быстрее находить эффективные стратегии маршрутизации и распределения потоков данных.

% Проверка указанных гипотез осуществляется путем аналитических расчетов и компьютерного моделирования, позволяя убедиться в правильности выдвинутых предположений и определить пределы их применимости.

% Работа построена на сочетании ряда общепринятых научных методов, среди которых выделяются:
% Общенаучные методы

%     Общелогические методы (анализ, абстрагирование, обобщение, идеализация, индукция, аналогия, моделирование и другие). Моделирование было ключевым инструментом для подтверждения гипотез и формирования общей картины функционирования сети.

% 	Конкретно-научные методы

% 	Особое значение имели специально-научные методы, присущие таким областям, как информатика и математика:
	
% 		Методы анализа графов. Применялись для оценки связности и устойчивости сетей.
% 		Вероятностные методы. Использовались для расчёта вероятности успешного приёма-передачи пакета данных.
% 		Имитационное моделирование. Позволило воспроизвести реалистичную картину функционирования сети и оценить предложенные решения.
	
% 	Применение вышеуказанной методологии обеспечило полноту охвата всех необходимых этапов исследования и позволило достичь поставленных целей и задач.


% Апробация результатов исследования выполнена в следующей форме:

%     Разработан конкретный проект автоматизированной системы передачи данных на основе предложенного алгоритма самоорганизации, который прошёл тестирование и показал высокие показатели отказоустойчивости и экономичности расхода энергии. Данный проект внедряется на территории тестового полигона крупной технологической компании и демонстрирует существенное улучшение показателей по сравнению с традиционными протоколами.
%     Основные итоги исследования закреплены в публикациях, которые вошли в список использованных источников. Среди них:
%         статья в международном журнале IEEE Transactions on Wireless Communications,
%         тезисы докладов на Всероссийской научно-технической конференции «Современные технологии телекоммуникаций»,
%         отчёт по результатам предварительного этапа исследований, размещённый на сайте факультета информационных технологий и программирования СПбПУ.

% Представленные публикации позволили привлечь дополнительное внимание специалистов и заинтересованных сторон к данному проекту, а также подтвердили оригинальность и новаторство предложенного подхода.

% Основная часть работы организована следующим образом:

%     Первая глава: Проводится подробный анализ и классификация самоорганизующихся низкоскоростных сетей передачи данных, рассматриваются их особенности, достоинства и ограничения. Изучаются физические и организационные аспекты построения таких сетей, предлагаются общие рекомендации по обеспечению их эффективной работы.
%     Вторая глава: Содержит разработку оригинального алгоритма функционирования самоорганизующейся сети, детально описывается архитектура и схема работы предложенного решения. Производятся расчеты и формируются математические модели, подтверждающие преимущество данного подхода.
%     Третья глава: Представляет реализацию прототипа сети на базе предложенного алгоритма, включая программное обеспечение и аппаратные компоненты. Подробно описаны этапы разработки, тестирование прототипа и получение первичных результатов. Показано, каким образом новая технология влияет на производительность и долговечность сети.
%     Четвёртая глава: Посвящена подведению итогов исследования, проведению серии контрольных измерений и сравнительного анализа полученной системы с аналогичными проектами. Сделаны выводы о преимуществах предложенного подхода и сформулированы дальнейшие перспективы развития данного направления.

% Такая организация материала способствует достижению главной цели исследования — разработке эффективного алгоритма функционирования самоорганизующейся низкоскоростной системы передачи данных.

% %% Вспомогательные команды - Additional commands
% %\newpage % принудительное начало с новой страницы, использовать только в конце раздела
% %\clearpage % осуществляется пакетом <<placeins>> в пределах секций
% %\newpage\leavevmode\thispagestyle{empty}\newpage % 100 % начало новой строки